% Direct-authority diplomatic French for printed p.133.
\SourcePage{133}

\LayerTarget{tohoku:example:1.7:h:p2}
\noindent diagrammes correspondants (pour une catégorie $\mathbf{C}$ donnée, non
nécessairement additive) sont connus sous le nom de \textit{systèmes
inductifs} sur $I$ à valeurs dans $\mathbf{C}$. Si on change $I$ en l'ensemble
préordonné opposé, ou $\mathbf{C}$ en $\mathbf{C}^\circ$, on obtient les \textit{systèmes
projectifs} sur $I$ à valeurs dans $\mathbf{C}$. Un cas important est celui où $I$ est
l'ensemble des parties ouvertes d'un espace topologique $X$, ordonné par la
relation $\supset$ : on obtient alors la notion de \textit{préfaisceau} sur
$X$ (à valeurs dans la catégorie $\mathbf{C}$).

\LayerTarget{tohoku:section:1.8}
\textbf{1.8. Limites inductives et projectives.} Nous ne parlerons
que des premières, la notion de limite projective étant duale de celle de
limite inductive. Soit $\mathbf{C}$ une catégorie, $I$ un ensemble préordonné,
$\mathbf A=(A_i,u_{ij})$ un système inductif sur $I$ à valeurs dans $\mathbf{C}$
($u_{ij}$ est un morphisme $A_j\to A_i$ défini pour $i\geqq j$).
\LayerTarget{tohoku:def:inductive-limit}
On appelle \textit{limite inductive} (généralisée) de $\mathbf A$ un système
formé d'un $\overline{\mathbf A}\in \mathbf{C}$ et d'une famille $(u_i)$ de morphismes
$u_i:A_i\to\mathbf A$, satisfaisant aux conditions suivantes :
a) Pour $i\leqq j$, on a $u_i=u_j u_{ji}$ b) Pour tout $B\in \mathbf{C}$ et toute
famille $(v_i)$ de morphismes $v_i:A_i\to B$, telle que l'on ait
pour tout couple $(i,j)$ avec $i\leqq j$, la relation $v_i=v_j u_{ji}$, on peut
trouver un morphisme $v$ et un seul de $\overline{\mathbf A}$ dans $B$ tel que l'on ait
$v_i=v u_i$ pour tout $i\in I$.
\LayerTarget{tohoku:def:inductive-limit:functoriality}
Si $(\overline{\mathbf A},(u_i))$ est une limite inductive de
$\mathbf A=(A_i,u_{ij})$, et si $(\mathbf B,(v_i))$ est une limite inductive
d'un deuxième système inductif $\mathbf B=(B_i,v_{ij})$, enfin si
$w=(w_i)$ est un morphisme de $\mathbf A$ dans $\mathbf B$, alors il existe un
morphisme $\overline w$ et un seul de $\overline{\mathbf A}$ dans $\overline{\mathbf B}$ tel que
l'on ait, pour tout $i\in I$ : $\overline w u_i=v_iw_i$. En particulier, deux
limites inductives d'un même système inductif sont canoniquement isomorphes
(dans un sens évident), aussi est-il naturel de choisir, pour tout système
inductif qui admet une limite inductive, une telle limite inductive (par
exemple au moyen du symbole $\tau$ de Hilbert), qu'on notera alors
$\mathop{\rm lim}\limits_{\longrightarrow}\mathbf A$ ou
$\mathop{\rm lim}\limits_{\substack{\longrightarrow\\i\in I}}A_i$, et qu'on
appellera \textit{la} limite inductive du système inductif donné.
\LayerTarget{tohoku:def:inductive-limit:functor}
Si $I$ et $\mathbf{C}$ sont tels que
$\mathop{\rm lim}\limits_{\longrightarrow}\mathbf A$ existe pour
\textit{tout} système inductif $\mathbf A$ sur $I$ à valeurs dans $\mathbf{C}$, il
résulte de ce qui précède que
$\mathop{\rm lim}\limits_{\longrightarrow}\mathbf A$ est un \textit{foncteur
covariant} défini sur la catégorie des systèmes inductifs de $I$ dans $\mathbf{C}$, à
valeurs dans $\mathbf{C}$.

\LayerTarget{tohoku:prop:1.8}
\textsc{Proposition 1.8.} \textit{Soit $\mathbf{C}$ une catégorie
abélienne satisfaisant à l'axiome \textup{AB 3)} (existence des sommes directes
quelconques) et soit $I$ un ensemble préordonné filtrant croissant. Alors pour
tout système inductif $\mathbf A$ sur $I$, à valeurs dans $\mathbf{C}$, la limite
inductive $\mathop{\rm lim}\limits_{\longrightarrow}\mathbf A$ existe, et
c'est un foncteur additif exact à droite de $\mathbf A$. Si $\mathbf{C}$ satisfait à
l'axiome \textup{AB 5)} \textup{(\textit{cf.} 1.5)} ce foncteur est même exact, et alors le noyau du
morphisme canonique
$u_i:A_i\to\mathop{\rm lim}\limits_{\longrightarrow}\mathbf A$
est le $\sup$ des noyaux des morphismes $u_{ji}:A_i\to A_j$ pour
$j\geqq i$. (En particulier, si les $u_{ji}$ sont injectifs, il en est de même
de $u_i$).}

\LayerTarget{tohoku:prop:1.8:proof:p1}
Pour construire une limite inductive de $(A_i,u_{ij})$, on considère $S=\displaystyle\bigoplus_{i\in I}A_i$, et pour tout couple $(i,j)$ avec $i\leqq j$, le morphisme $w_{ij}$ de
$A_i$ dans $S$ déduit du morphisme $A_i\to A_i+A_j$ dont les
composantes sont $I_{A_i}$ et $-u_{ji}$. Soit $N_{ij}$

\TohokuContinuation{134}
